% Source presentations for register row N-c-025.

\subsection{The finite-template CSP dichotomy: problem, conjecture, and two proofs}
\label{sec:p09-presentation-n-c-025}

\paragraph{The fixed-template problem.}
Fix a finite relational signature \(\tau\): a finite list of relation symbols,
each with a specified positive arity.  A finite \(\tau\)-structure \(A\)
consists of a finite domain \(\lvert A\rvert\) and, for each \(r\)-ary
\(R\in\tau\), a relation \(R^A\subseteq\lvert A\rvert^r\).  For two such
structures, a \emph{homomorphism} \(h:A\to B\) is a function
\(h:\lvert A\rvert\to\lvert B\rvert\) such that
\[
 (a_1,\ldots,a_r)\in R^A
 \quad\Longrightarrow\quad
 (h(a_1),\ldots,h(a_r))\in R^B
\]
for every \(R\in\tau\) and every tuple in \(R^A\).  Preservation, rather than
reflection, of the displayed relations is the required distinction from an
embedding or an isomorphism.  The definition is used immediately: for a fixed
finite template \(B\),
\[
 \operatorname{CSP}(B)
   =\{A:A\text{ is a finite }\tau\text{-structure and }A\to B\}.
\]
Only \(A\) is input.  If \(B\) were also input, the result would be the
different, uniform finite-homomorphism problem
\cite[\S\S2--3]{FederVardi1998}.

The constraint and homomorphism presentations are literally the same problem.
Let \(D\) be finite and let
\(\Gamma=\{R_1,\ldots,R_s\}\) be a finite set of finitary relations on
\(D\), called a \emph{finite constraint language}.  A
\(\operatorname{CSP}(\Gamma)\) instance has a finite variable set \(V\) and a
finite conjunction of constraints
\[
 R_i(v_1,\ldots,v_{r_i}),
 \qquad R_i\in\Gamma,
\]
and asks for an assignment \(g:V\to D\) whose value tuple belongs to \(R_i\)
at every displayed constraint.  Give the template \(B_\Gamma\) domain \(D\)
and relation \(R_i^{B_\Gamma}=R_i\).  Give the input structure \(A_I\) domain
\(V\) and place a tuple in \(R_i^{A_I}\) exactly when that tuple occurs as an
\(R_i\)-constraint of \(I\).  Then
\[
 g\text{ satisfies }I
 \quad\Longleftrightarrow\quad
 g:A_I\to B_\Gamma\text{ is a homomorphism}.
\]
Thus a constraint language names the fixed allowed relations, whereas an
instance names the varying variables and applications of those relations.
Neither is a second kind of solution certificate.

Here is the binary case in both notations.  Let \(\tau=\{E\}\) and fix
\[
 \lvert B_2\rvert=\{0,1\},
 \qquad
 E^{B_2}=\{(0,1),(1,0)\}.
\]
Let \(A_P\) have domain \(\{a,b,c\}\) and
\[
 E^{A_P}=\{(a,b),(b,a),(b,c),(c,b)\}.
\]
In constraint notation the instance says
\(x_a\ne x_b\) and \(x_b\ne x_c\).  The assignment
\[
 h(a)=0,\qquad h(b)=1,\qquad h(c)=0
\]
is a solution precisely because it is a homomorphism \(A_P\to B_2\).  Adding
both orientations of the edge \(\{a,c\}\) produces a triangle and makes the
instance negative: the first two edges force \(a\) and \(c\) to have the same
value, while the third requires different values.  More generally,
\(\operatorname{CSP}(K_k)\) is graph \(k\)-colorability, with \(K_k\) fixed
and the graph to be colored varying.

\paragraph{The logical class and the exact reduction types.}
The name \emph{monotone monadic SNP without inequality}---where \emph{SNP}
abbreviates \emph{strict NP}---is needed to distinguish a syntactically
specified class of finite-structure decision problems from the fixed-template
problems just defined.  Fix a finite input signature \(\tau\).  An MMSNP
sentence has an equality-free normal form
\[
 \Phi=
 \exists M_1\cdots\exists M_q\,
 \forall x_1\cdots\forall x_n\,
 \bigwedge_{i=1}^{m}
   \neg\bigl(\alpha_i(\bar x)\land\beta_i(\bar x)\bigr),
\]
where every \(M_j\) is a unary relation symbol not in \(\tau\), every
\(\alpha_i\) is a conjunction of positive atomic formulas over \(\tau\), and
every \(\beta_i\) is a conjunction of literals \(M_j(x)\) or
\(\neg M_j(x)\).  The tuple of existential symbols may be empty, and neither
equality nor inequality occurs in the displayed normal form.  Thus
\emph{monadic} restricts the existential second-order relations to unary
relations, \emph{strict} leaves only a universal first-order block after that
existential block, and \emph{monotone} says that each input relation occurs
negatively in the whole sentence.  A fixed sentence \(\Phi\) defines the
problem
\[
 L_\Phi=\{A:A\text{ is a finite }\tau\text{-structure and }A\models\Phi\}.
\]
MMSNP therefore stands here for a sentence fragment and the associated class
of problems; it is not a new kind of structure, template, or homomorphism.

Every fixed-template problem \(\operatorname{CSP}(B)\) is definable in this
fragment.  Introduce one unary predicate \(M_b\) for each
\(b\in\lvert B\rvert\), require every input element to have exactly one such
color, and, for each relation \(R^B\), forbid precisely the colored input tuples
whose color tuple does not belong to \(R^B\).  A coloring satisfying those
clauses is exactly a homomorphism into \(B\).  The converse inclusion is not
literal.  On finite simple graphs, for example,
\[
 \forall x\,\forall y\,\forall z\,
 \neg\bigl(E(x,y)\land E(y,z)\land E(z,x)\bigr)
\]
is an MMSNP sentence defining triangle-freeness, but Feder and Vardi show that
this language is not \(\operatorname{CSP}(B)\) for any fixed finite graph
\(B\).  This witnesses the nonredundancy of the logical class even though the
two classes have the following computational relationship.

The deterministic use of \emph{Karp reduction} below is the one-output
polynomial-time many-one map fixed in
Convention~\ref{def:p04-complexity-convention}.  By contrast, the stated
Turing reduction may consult the target decision problem on more than one
instance; \emph{randomized} records that the reducing computation itself uses
random choices.  These are reduction types, not additional problem classes.

Feder--Vardi's reduction theorem, Theorem~6 in the complete author-hosted
manuscript, states that for every \(L_\Phi\) in MMSNP there is a
finite-template CSP problem \(\Omega_\Phi\) such that \(L_\Phi\) has a
deterministic polynomial-time Karp reduction to \(\Omega_\Phi\), whereas
\(\Omega_\Phi\) has a randomized polynomial-time Turing reduction to
\(L_\Phi\).  The source further says that the latter direction can be a
\emph{randomized} Karp reduction when the CSP input's Gaifman graph, as
defined in Definition~\ref{def:p09-locality}, is connected; the general
disconnected input is handled component by component
\cite[author-copy pp.~12--13, thm.~6]{FederVardi1998}.

\paragraph{The randomized relational-girth construction.}
The proof needs a measure of cycles that treats an occurrence of an
\(r\)-ary relation as one \(r\)-place constraint.  For a finite relational
structure \(S\), say that \(S\) has \emph{relational girth greater than
\(k\)} when every family of \(m\) relation occurrences, where
\(1\leq m\leq k\) and the \(j\)-th occurrence has arity \(r_j\),
mentions at least
\[
  1+\sum_{j=1}^{m}(r_j-1)
\]
distinct elements.  This is a property of a relational structure, not of a
formula or a reduction.  It is not the girth of the Gaifman graph: one
high-arity tuple already induces a clique in that graph, while it counts as
one relation occurrence here.  The separate term is therefore necessary for
the arity-sensitive acyclicity condition used in the next construction and in
the subsequent forbidden-pattern argument.

Feder--Vardi Theorem~5 fixes positive integers \(k\) and \(d\).  For every
\(n\)-element structure \(S\), it constructs in randomized polynomial time
a structure \(S^{\ast}\), of size \(n^a\) for a constant \(a\) depending
only on \(k\) and \(d\), such that \(S^{\ast}\) has relational girth
greater than \(k\) and \(S^{\ast}\to S\), and, for every structure \(T\)
over the same signature with at most
\(d\) elements,
\[
 S^{\ast}\to T \quad\Longleftrightarrow\quad S\to T.
\]
Here \(S^{\ast}\) denotes exactly the finite structure produced by this
construction; it is not a new kind of CSP instance
\cite[author-copy pp.~10--11, thm.~5]{FederVardi1998}.

To construct it, take \(N=n^s\) copies of every element of \(S\), where
\(s\) is a sufficiently large constant.  For each original occurrence of
an \(r\)-ary relation and each choice of copies of its entries, independently
retain the lifted occurrence with probability
\[
 p_r=N^{1-r+\varepsilon},
\]
where \(\varepsilon>0\) is a sufficiently small constant; unary occurrences
are put on every copy.  Thus an original \(r\)-ary occurrence has expected
\(N^r p_r=N^{1+\varepsilon}\) retained lifts.

Now consider a prospective short relational cycle containing \(m\leq k\)
occurrences of arities \(r_1,\ldots,r_m\) and mentioning \(t\) copied
elements.  Its failure of the displayed girth condition gives
\(t\leq\sum_j(r_j-1)\).  For a fixed underlying pattern, the expected
number of copied realizations is at most
\[
 N^t\prod_{j=1}^{m}N^{1-r_j+\varepsilon}
 =N^{t+m-\sum_jr_j+m\varepsilon}
 \leq N^{k\varepsilon}.
\]
Only \(n^{a_0}\) underlying patterns need be considered for a constant
\(a_0\), so the expected number of short cycles is at most
\(N^\delta\), where \(\delta=a_0/s+k\varepsilon\).  Markov's inequality
therefore gives, with probability at least one half, at most
\(2N^\delta\) such cycles.  Delete one relation occurrence from each of
them.  The resulting structure has relational girth greater than \(k\),
and projection of every copy to its original element is the homomorphism
\(S^{\ast}\to S\).  Hence \(S\to T\) always implies
\(S^{\ast}\to T\).

For the converse, suppose \(h:S^{\ast}\to T\) and
\(\lvert T\rvert\leq d\).  Among the \(N\) copies of each original
element \(u\), at least \(N/d\) receive one common value under \(h\);
choose \(N/d\) such copies.  For every original \(r\)-ary relation
occurrence, the expected number of retained lifts using prescribed choices of
these subsets is
\[
 (N/d)^r p_r=N^{1+\varepsilon}/d^r.
\]
A Cherno lower-tail bound makes the probability of retaining fewer than half
this number exponentially small in \(N^{1+\varepsilon}\).  There are at
most polynomially many original relation occurrences and at most
\(2^{rN}\) choices of the relevant subsets for each one.  A union bound
therefore shows that, with positive probability, every such choice has the
required supply of lifts simultaneously.  Choose \(s\) large and
\(\varepsilon\) small so that \(\delta<1+\varepsilon\); deleting at
most \(2N^\delta\) short-cycle occurrences then cannot remove all the
lifts for any one choice.  Assign to \(u\) the common \(h\)-value of
its chosen copies.  A surviving lift of every original relation occurrence
shows that this assignment preserves that relation.  It is therefore a
homomorphism \(S\to T\), proving the equivalence.

\paragraph{The complete author-copy reduction.}
Theorem~6 begins with an MMSNP sentence written as a conjunction of forbidden
conjunctions.  A forbidden conjunction is called \emph{biconnected} here
when it cannot be split either into two parts with disjoint variable sets or
into two parts whose variable sets share only one variable.  This is a
syntactic property of one forbidden conjunction.  It is stronger than mere
connectedness because it excludes a single articulation variable, and it is
needed below to force a forbidden pattern in a high-girth structure to lie
inside one explicit relation occurrence.

The source reduces every forbidden conjunction to that form.  If
\(A(\bar x)\) and \(B(\bar y)\) have disjoint variables, then
\[
 \neg(A(\bar x)\land B(\bar y))
\]
is replaced, under a fresh existential zero-ary symbol \(p\), by
\[
 (A(\bar x)\to p)\land(B(\bar y)\to\neg p).
\]
A zero-ary existential relation symbol is simply a globally chosen truth
value, not a predicate on domain elements.  If the two parts share exactly
one variable \(z\), a fresh existential unary predicate \(q\) instead
gives
\[
 (A(\bar x,z)\to q(z))
 \land(B(\bar y,z)\to\neg q(z)).
\]
Before these decompositions, the formula is closed under all identifications
of variables occurring in a forbidden conjunction, since identifying two
variables can destroy biconnectedness.  If an input atom has all arguments
equal and occurs with other input atoms, it is replaced by a fresh monadic
proxy at that element, together with the implication from the original
repeated-argument atom to the proxy.  These preparations account for the
degenerate identifications rather than assuming them away.

For each remaining forbidden conjunction, let \(C\) be the conjunction of
its input-signature atoms.  Introduce one fresh input relation \(R_C\), whose
arguments are all the distinct variables of the entire forbidden conjunction,
and replace \(C\) by \(R_C\).  Whenever the input part \(C'\) of another forbidden
conjunction is a sub-conjunction of \(C\), also add the corresponding
instance obtained by extending \(C'\) to \(C\), using fresh variables
where necessary.  The resulting fixed formula has at most one input relation
in each forbidden conjunction, and that relation is applied to distinct
variables.  From an original input structure, its expanded structure is
computed deterministically by declaring \(R_C\) on an assignment to all those
variables exactly when the input atoms in \(C\) hold under that assignment.

The reverse substitution \(R_C\mapsto C\) can create an occurrence of
\(R_C\) that was not explicit in the expanded input.  The source's example
replaces a triangle by a ternary relation: three explicit ternary tuples can,
after expansion back to edges, assemble a fourth triangle.  Choose \(k\)
larger than the number of input atoms in any \(C\), and apply Theorem~5.
In the resulting high-relational-girth instance, a biconnected forbidden
conjunction assembled from at most \(k\) relation occurrences must be
contained in one of them.  The sub-conjunction closure just imposed then says
that the corresponding prohibition was already present for that explicit
occurrence.  Thus high relational girth removes precisely the obstruction to
the reverse substitution.

It remains to recognize the special formula as a finite-template CSP.  First
suppose there are no zero-ary existential symbols.  The template has one
element for every allowed complete truth assignment to the existential unary
predicates at a single domain element.  For each input relation, its template
relation contains exactly those tuples of such unary types that the forbidden
conjunctions do not exclude.  A homomorphism into this template is therefore
the same thing as a choice of the existential unary predicates satisfying the
special formula.  If zero-ary existential symbols occur, make one such finite
template \(T_i\) for each truth assignment to them, take the disjoint union
of the \(T_i\), and add a binary relation holding precisely on pairs of
elements from the same \(T_i\).  A connected input can map into only one
\(T_i\); disconnected inputs are decided component by component.

The deterministic expansion by the relations \(R_C\) is the Karp reduction
from MMSNP to CSP.  In the reverse direction, Theorem~5 supplies the randomized
high-girth replacement before the relations \(R_C\) are expanded, giving a
randomized Karp reduction on connected inputs.  Asking separately about every
connected component gives the randomized Turing reduction on arbitrary
inputs.  This reconstructs the source's reduction orientation rather than
inferring it from the phrase ``polynomially equivalent''
\cite[author-copy pp.~12--13, thm.~6]{FederVardi1998}.

\paragraph{Source boundary.}
The proof just reconstructed was inspected in the complete fifty-three-page
PostScript manuscript linked from Tom\'as Feder's Stanford publications page;
Theorems~5 and 6 occupy author-copy pages 10--11 and 12--13, respectively.
That file is an author-hosted complete manuscript, not a facsimile of the
SIAM typesetting.  The journal identity and DOI are independently fixed by
the publisher record, but the page locators in the preceding paragraphs are
author-copy locators.

The proof of Theorem~5 explicitly left derandomization open.  Kun later
replaced the randomized sparse-incomparability step by a
deterministic polynomial-time construction.  His statement uses a finite
family \(\mathcal T\) of templates and defines
\[
 \operatorname{CSP}(\mathcal T)
   =\bigcup_{B\in\mathcal T}\operatorname{CSP}(B).
\]
This notation stands for a finite union of fixed-template problem languages,
not for a second kind of input instance.  Kun gives deterministic
polynomial-time reductions in both directions between \(L_\Phi\) and such a
\(\operatorname{CSP}(\mathcal T)\), thereby removing the random step without
retroactively changing the 1998 theorem \cite{Kun2013}.
Computational equivalence here is neither literal equality of the two problem
classes nor a proof of a complexity dichotomy.
\index{MMSNP}\index{constraint-satisfaction problem!MMSNP correspondence}
\index{constraint-satisfaction problem!finite family of templates}

\paragraph{Conjecture in 1998, theorem in 2017.}
The Feder--Vardi journal article published in 1998 organizes this fixed-template
framework, connects it to logical and algebraic methods, and formulates a
dichotomy program.  The reduction theorem above transfers the classification
question between MMSNP and finite-template CSPs, but the article's abstract
expressly says that the desired dichotomy is not achieved.  It is therefore a
source for the problem, reduction, and conjecture, not a proof of
Theorem~\ref{thm:p09-csp-dichotomy}.  The status changed at FOCS 2017, where
Bulatov and Zhuk gave independent proofs of the finite-domain conjecture
\cite{FederVardi1998,Bulatov2017,Zhuk2017}.

\paragraph{The normalizations required by the algebraic criterion.}
Let \(\Gamma\) be a finite constraint language on \(D\).  An
\emph{endomorphism of \(\Gamma\)} is a unary function \(e:D\to D\) preserving
every relation in \(\Gamma\); equivalently, it is a unary polymorphism in the
definition below.  The language is a \emph{core} when every endomorphism is a
bijection.  Corehood names a property of the fixed language, not of an input
instance, and rules out a proper endomorphic image that represents the same
constraint problem.

For an \(r\)-ary relation \(R\), define
\[
 e(R)=\{(e(a_1),\ldots,e(a_r)):(a_1,\ldots,a_r)\in R\},
 \qquad
 e(\Gamma)=\{e(R):R\in\Gamma\},
\]
with domain \(e(D)\).  Choose an endomorphism \(e\) whose image has minimum
cardinality among endomorphic images of \(\Gamma\).  Zhuk's Theorem~1 states
that \(e(\Gamma)\) is a core and that
\(\operatorname{CSP}(\Gamma)\) and
\(\operatorname{CSP}(e(\Gamma))\) reduce to one another in polynomial time.
This is the \emph{core reduction}; the name abbreviates exactly this passage to
a minimum endomorphic image and is used below to transport the dichotomy
criterion back to the original language.

For a core language on \(D\), put
\[
 \Gamma^+=\Gamma\cup\bigl\{\{d\}:d\in D\bigr\}.
\]
Zhuk's Theorem~2 gives a polynomial-time reduction from
\(\operatorname{CSP}(\Gamma^+)\) to \(\operatorname{CSP}(\Gamma)\); the
reverse reduction merely regards a \(\Gamma\)-instance as a
\(\Gamma^+\)-instance.  The operation is the \emph{singleton expansion}: it
adds the unary relation fixing each domain element.  It differs from core
reduction because it does not shrink the domain, and it is needed to enforce
idempotence of the operations in the criterion
\cite[pp.~331--332, thms.~1--2]{Zhuk2017}.

The idempotence claim is immediate but worth displaying.  If an \(m\)-ary
operation \(f\) preserves every singleton \(\{d\}\), then applying \(f\) to
\(m\) copies of the unary tuple \((d)\) gives
\[
 f(d,\ldots,d)\in\{d\},
 \qquad\text{hence}\qquad
 f(d,\ldots,d)=d.
\]
Conversely, an idempotent operation preserving \(\Gamma\) automatically
preserves all the added singleton relations.  Thus the core and singleton
steps are polynomially equivalent changes to the fixed language; they are not
promises silently imposed on each input.

\paragraph{Polymorphisms and the WNU boundary.}
Let \(m\geq1\).  An operation \(f:D^m\to D\) is a
\emph{polymorphism of \(\Gamma\)} if it preserves every relation in
\(\Gamma\) coordinatewise.  Explicitly, whenever \(R\in\Gamma\) is
\(r\)-ary and
\[
 \mathbf a^j=(a^j_1,\ldots,a^j_r)\in R
 \qquad(1\leq j\leq m),
\]
one has
\[
 \bigl(f(a^1_1,\ldots,a^m_1),\ldots,
        f(a^1_r,\ldots,a^m_r)\bigr)\in R.
\]
The operation acts on the fixed domain and uniformly preserves every allowed
relation.  This standing distinguishes it from a constraint, which belongs to
one input, and from a satisfying assignment, which maps the variables of one
input into \(D\).  The distinction is necessary because the dichotomy
criterion concerns operations shared by all instances over the template.

For \(m\geq3\), a polymorphism \(w:D^m\to D\) is a
\emph{weak near-unanimity operation}, abbreviated \emph{WNU}, when, for every
\(x,y\in D\),
\[
 w(x,\ldots,x)=x
\]
and
\[
 \begin{split}
 w(y,x,\ldots,x)
   &=w(x,y,x,\ldots,x)\\
   &=\cdots
    =w(x,\ldots,x,y).
 \end{split}
\]
The common value in the second display need not be \(x\); requiring that
additional equality would give the stronger near-unanimity equations.  A WNU
therefore stands for a polymorphism satisfying exactly the two displayed
conditions, is not merely an arbitrary polymorphism, and is the operation used
in the following normalized criterion.

Start with \(\Gamma\), pass to a minimum endomorphic core \(\Gamma_c\), and
form its singleton expansion \(\Gamma_c^+\).  The finite-domain criterion is
then:
\[
 \begin{array}{c|c}
 \Gamma_c^+\text{ has a WNU polymorphism}
   & \operatorname{CSP}(\Gamma)\in\mathsf P,\\[2pt]
 \Gamma_c^+\text{ has no WNU polymorphism}
   & \operatorname{CSP}(\Gamma)\text{ is }\mathsf{NP}\text{-complete}.
 \end{array}
\]
The polynomial equivalences above are what make the criterion a classification
of the original \(\Gamma\), not only of its normalized presentation.

\paragraph{What the two 2017 papers prove.}
Bulatov's Theorem~1 states the finite nonuniform dichotomy: every finite
constraint language over a finite set gives a CSP in \(\mathsf P\) or an
\(\mathsf{NP}\)-complete CSP.  His Theorem~4 gives, for the associated finite
idempotent setting, the WNU characterization used in the normalized statement
above.  The implication from failure of the algebraic condition to hardness
came from earlier algebraic CSP work.  Bulatov's new direction is the
polynomial-time algorithm for the WNU side.  The FOCS paper states the
structural and algorithmic spine; the author-hosted expanded manuscript gives
the long proof \cite[thms.~1 and 4]{Bulatov2017}.

Zhuk independently proves the same finite nonuniform dichotomy.  His
Theorems~1--2 are the core and singleton reductions just stated, and
Theorem~3 records the previously known hard direction when the normalized
language has no WNU polymorphism.  His new contribution is the polynomial-time
algorithm when such a polymorphism exists, together with its correctness
argument.  The conference paper supplies the principal reductions and proof
architecture and sends omitted auxiliary proofs to the full version
\cite[pp.~331--336 and 341--342]{Zhuk2017}.

Neither proof is presented here as a consequence of the other, and neither
algorithm is replaced by the WNU equations alone.  Those algorithms depend on
long structural analyses that lie beyond this source presentation's declared
citation boundary.  What has been reconstructed locally is the exact problem,
the two normalizations, the algebraic criterion, the historical division
between prior hardness and the new tractable algorithms, and the scope of the
cited theorem.

Theorem~\ref{thm:p09-csp-dichotomy} classifies finite-domain,
fixed-template, decision CSPs.  It does not classify counting or optimization
CSPs, promise CSPs, problems with a varying template, or arbitrary
infinite-domain CSPs.  Changing any of those items changes the problem rather
than adding a harmless notational variant.
